\documentclass[12pt, a4paper]{article}
\usepackage{amsmath, amssymb, amsthm}
\usepackage{geometry}
\usepackage{enumitem}
\usepackage{xcolor}
\usepackage{mdframed}
\usepackage{hyperref}
\usepackage{booktabs}
\usepackage{ctex}
\usepackage{bm}
\usepackage{listings}
\geometry{margin=2.5cm}

\newcounter{probcnt}
\setcounter{probcnt}{0}
\newenvironment{problem}[1][]{%
  \stepcounter{probcnt}
  \begin{mdframed}[linecolor=black!30, backgroundcolor=gray!4,
    innertopmargin=8pt, innerbottommargin=8pt,
    innerleftmargin=10pt, innerrightmargin=10pt,
    skipabove=10pt, skipbelow=4pt]
  \noindent\textbf{题 \theprobcnt}\ifx&#1&\else\ \textbf{[#1]}\fi\quad
}{%
  \end{mdframed}
}
\newenvironment{subproblems}{%
  \begin{enumerate}[label=(\alph*), leftmargin=1.8em, topsep=4pt, itemsep=4pt]
}{%
  \end{enumerate}
}
\newcommand{\hint}[1]{\smallskip\noindent\textit{提示：#1}}
\newcommand{\context}[1]{\smallskip\noindent\textcolor{gray}{\small \textbf{背景：}#1}\smallskip}

\definecolor{myblue}{RGB}{0, 102, 204}
\definecolor{mygreen}{RGB}{0, 153, 76}
\lstset{
  language=Python,
  basicstyle=\ttfamily\small,
  keywordstyle=\color{myblue}\bfseries,
  commentstyle=\color{gray},
  stringstyle=\color{mygreen},
  numbers=left, numberstyle=\tiny\color{gray},
  breaklines=true, frame=single,
  backgroundcolor=\color{gray!5}
}

\newcommand{\R}{\mathbb{R}}

\title{\textbf{Week 3 作业：算法题、LP 松弛与 GPU}}
\author{优化算法课程}
\date{}

\begin{document}
\maketitle
\thispagestyle{empty}

\noindent\textbf{说明}：这周的核心想法是：
\begin{itemize}[leftmargin=1.8em]
    \item 很多熟悉的算法题都能写成整数规划（ILP）；
    \item 把整数约束放松后，可以得到线性规划（LP）；
    \item LP 松弛后可能不再精确，但形式更规整，往往更适合 GPU 并行。
\end{itemize}

\noindent 请尽量用自己的话解释：\textbf{原问题是什么、LP 松弛后变成了什么、误差从哪里来、为什么更适合 GPU。}



\bigskip

\begin{problem}[基础]
\textbf{一元线性回归中的梯度下降步长}

\context{考虑一元线性回归模型
$$
\hat y_i = wx_i,
$$
其中参数 $w\in\R$，训练数据为 $(x_i,y_i)$，$i=1,\dots,n$。平方损失函数定义为
$$
f(w)=\frac{1}{2n}\sum_{i=1}^n (wx_i-y_i)^2.
$$
我们用梯度下降更新参数：
$$
w_{t+1}=w_t-\eta f'(w_t).
$$}

\begin{subproblems}
    \item 计算目标函数 $f(w)$ 的一阶导数 $f'(w)$ 和二阶导数 $f''(w)$。
    
    \item 设
    $$
    A=\frac{1}{n}\sum_{i=1}^n x_i^2.
    $$
    证明梯度下降更新可以写成
    $$
    w_{t+1}-w^\ast = (1-\eta A)(w_t-w^\ast),
    $$
    其中 $w^\ast$ 是最优解。
    
    \item 由上式说明：为了使梯度下降收敛，步长 $\eta$ 至少应满足什么条件？请给出 $\eta$ 的最大允许范围。
    
    \item 用自己的话解释：为什么当数据的 $x_i$ 整体更大时，允许的最大步长会更小？这说明了步长与数据尺度之间有什么关系？
    
    \item （思考）若把所有输入同时放大为 $x_i' = c x_i$，其中 $c>0$，那么允许的最大步长会如何变化？
\end{subproblems}
\end{problem}


% -------------------------------------------------------
\begin{problem}[基础]
\textbf{最大权二分匹配}

\context{有 $n$ 个用户、$m$ 个广告位。每个广告位最多分给一个用户，每个用户最多得到一个广告位。兼容度矩阵为 $C\in\R^{n\times m}$。平台希望最大化总兼容度。}

\begin{subproblems}
    \item 用变量 $x_{ij}$ 表示“用户 $i$ 是否分配给广告位 $j$”，写出该问题的 ILP。
    \item 去掉整数约束，写出 LP 松弛。
    \item 用一句话说明：为什么这个问题的 LP 松弛比较“幸运”，通常不会带来误差？
\end{subproblems}
\end{problem}

% -------------------------------------------------------
\begin{problem}[基础]
\textbf{0-1 背包}

\context{有 $n$ 个物品，第 $i$ 个物品价值为 $v_i$，重量为 $w_i$，背包容量为 $W$。每个物品只能选或不选。}

\begin{subproblems}
    \item 写出 0-1 背包的 ILP。
    \item 写出对应的 LP 松弛。
    \item 解释为什么 LP 松弛后可能出现“拿半个物品”的现象，因此 LP 最优值一般会比整数最优值更乐观。
\end{subproblems}
\end{problem}

% -------------------------------------------------------
\begin{problem}[基础]
\textbf{Vertex Cover}

\context{给定无向图 $G=(V,E)$。希望选出尽可能少的顶点，使得每条边至少有一个端点被选中。}

\begin{subproblems}
    \item 写出 Vertex Cover 的 ILP。
    \item 写出对应的 LP 松弛。
    \item 若 LP 解中有些顶点取值是 $1/2$，这说明了什么？请用一句话解释“松弛误差”的含义。
\end{subproblems}
\end{problem}

\bigskip

% -------------------------------------------------------
\begin{problem}[理解]
\textbf{背包问题中的松弛误差}


\context{考虑如下 0-1 背包问题：共有 3 个物品，
$$
(v_1,w_1)=(14,6),\qquad (v_2,w_2)=(12,5),\qquad (v_3,w_3)=(8,4),
$$
背包容量为 $W=10$。设决策变量 $x_i\in\{0,1\}$ 表示是否选择物品 $i$。}

\begin{subproblems}
    \item 写出该问题的 ILP。
    
    \item 忽略整数约束，写出对应的 LP 松弛，并用“允许最后一个物品取分数”的方式，计算根节点的 LP 上界。
    
    \item 以变量 $x_1$ 为第一个分支变量，分别考虑两个子问题：
    \begin{enumerate}[label=\roman*., leftmargin=2em]
        \item $x_1=1$；
        \item $x_1=0$。
    \end{enumerate}
    对每个子问题，计算其 LP 松弛上界，并判断该分支是否可能产生更优整数解。
    
    \item 根据上面的结果，继续做一步分支定界，找出该背包问题的最优整数解。
    
    \item 用自己的话解释：分支定界法中的“分支”和“定界”分别是什么意思？为什么 LP 松弛可以用来帮助剪枝？
\end{subproblems}
\end{problem}

% -------------------------------------------------------
\begin{problem}[理解]
\textbf{Vertex Cover 的简单舍入}

\context{对 Vertex Cover 的 LP 松弛，考虑如下简单规则：若 $x_v^\ast \ge 1/2$，就把顶点 $v$ 选入解中。}

\begin{subproblems}
    \item 说明为什么这个规则一定得到一个合法的顶点覆盖。
    \item 证明该规则得到的解大小不超过 $2\sum_v x_v^\ast$。
    \item 用自己的话解释：为什么这说明 LP 松弛虽然不精确，但仍然能给出一个“还不错”的近似解？
\end{subproblems}
\end{problem}

% -------------------------------------------------------
\begin{problem}[理解]
\textbf{最短路为什么几乎没有松弛误差}

\context{在有向图中，最短路可以写成
$$
\max d_t
\qquad
\text{s.t.}\quad d_v-d_u\le w_{(u,v)},\quad d_s=0.
$$}

\begin{subproblems}
    \item 这里的变量 $d_v$ 表示什么？
    \item 为什么这个模型本身已经是连续的，不像背包或 Vertex Cover 那样需要“选或不选”？
    \item 用一句话比较：为什么有些问题松弛后误差大，而最短路这种问题几乎没有这个问题？
\end{subproblems}
\end{problem}

\bigskip

% -------------------------------------------------------
\begin{problem}[讨论]
\textbf{从离散算法到规整迭代}

\context{很多 CPU 上的经典算法（如 Dijkstra、动态规划、分支搜索）都依赖复杂控制流；而 GPU 更擅长做大量结构相同的并行计算。}

\begin{subproblems}
    \item 请比较下面两类计算的区别：
    \begin{enumerate}[label=\roman*., leftmargin=2em]
        \item 用堆实现 Dijkstra；
        \item 对所有边同时做一次 Bellman-Ford 松弛。
    \end{enumerate}
    哪一种更适合 GPU？为什么？
    
    \item 请比较下面两类计算的区别：
    \begin{enumerate}[label=\roman*., leftmargin=2em]
        \item 用回溯法精确求解 0-1 背包：不断尝试“选 / 不选”每个物品；
        \item 将 0-1 背包放松成连续变量问题后，同时计算所有物品变量的更新方向。
    \end{enumerate}
    哪一种更适合 GPU？为什么？
    \end{subproblems}
\end{problem}

% -------------------------------------------------------
\begin{problem}[讨论]
\textbf{编辑距离与反对角线并行}

\context{编辑距离的经典动态规划为
$$
\mathrm{dp}[i][j]
=
\min\Bigl(
\mathrm{dp}[i-1][j]+1,\ 
\mathrm{dp}[i][j-1]+1,\ 
\mathrm{dp}[i-1][j-1]+\mathbf 1[s_i\neq t_j]
\Bigr).
$$}

\begin{subproblems}
    \item 解释为什么同一条反对角线上的状态可以同时更新。
    \item 说明这类“分层更新”的结构为什么适合 GPU。
    \item 用一句话说明：这道题虽然本质上是 DP，但为什么也能看成“更规整的并行松弛过程”？
\end{subproblems}
\end{problem}

\bigskip

% -------------------------------------------------------
\begin{problem}[编程]
\textbf{值迭代：CPU 与 GPU 对比}

\context{值迭代的核心更新为
$$
V(s)\leftarrow \max_a \left[r(s,a)+\gamma\sum_{s'}P(s'|s,a)V(s')\right].
$$
当对所有状态同时更新时，这一过程很适合向量化实现。}

\textbf{任务}：在一个网格 MDP 上，实现 CPU 版和 GPU 版值迭代，并比较运行时间。

\begin{subproblems}
    \item 实现一个网格 MDP 生成函数，返回奖励矩阵 $R$ 和转移张量 $P$。
    \item 用 NumPy 实现 CPU 版值迭代。
    \item 用 PyTorch 实现 GPU 版值迭代，并正确计时。
    \item 在不同网格规模下比较 CPU 与 GPU 的运行时间，并回答：
    \begin{enumerate}[label=\roman*., leftmargin=2em]
        \item 状态数变大后，GPU 何时开始占优？
        \item 每轮迭代中，最耗时的部分通常是什么？
        \item 为什么值迭代比很多“带大量分支的搜索算法”更适合 GPU？
    \end{enumerate}
\end{subproblems}

\textbf{参考代码框架如下：}

\begin{lstlisting}
import numpy as np
import torch
import time

def make_grid_mdp(grid_size, gamma=0.99, seed=42):
    """
    生成 grid_size x grid_size 的网格 MDP。
    返回：
        R: shape (S, A)
        P: shape (A, S, S)
    """
    pass

def value_iteration_cpu(R, P, gamma=0.99, tol=1e-6, max_iter=10000):
    S, A = R.shape
    V = np.zeros(S)
    start = time.time()
    for it in range(max_iter):
        Q = R + gamma * np.einsum('aij,j->ia', P, V)
        V_new = Q.max(axis=1)
        if np.max(np.abs(V_new - V)) < tol:
            return V_new, it + 1, time.time() - start
        V = V_new
    return V, max_iter, time.time() - start

def value_iteration_gpu(R_np, P_np, gamma=0.99, tol=1e-6,
                        max_iter=10000, device='cuda'):
    R = torch.tensor(R_np, dtype=torch.float32, device=device)
    P = torch.tensor(P_np, dtype=torch.float32, device=device)
    V = torch.zeros(R_np.shape[0], dtype=torch.float32, device=device)

    torch.cuda.synchronize()
    start = time.time()

    for it in range(max_iter):
        Q = R + gamma * torch.einsum('aij,j->ia', P, V)
        V_new = Q.max(dim=1).values
        if torch.max(torch.abs(V_new - V)).item() < tol:
            torch.cuda.synchronize()
            return V_new.cpu().numpy(), it + 1, time.time() - start
        V = V_new

    torch.cuda.synchronize()
    return V.cpu().numpy(), max_iter, time.time() - start
\end{lstlisting}

\end{problem}

\end{document}